\documentclass[11pt,a4paper]{article}

% ---------- Layout ----------
\usepackage[margin=1in]{geometry}
\usepackage{microtype}

% ---------- Math ----------
\usepackage{amsmath,amssymb,amsthm,mathtools}

% ---------- Tables, links, bibliography ----------
\usepackage{booktabs}
\usepackage{graphicx}
\usepackage{hyperref}
\hypersetup{colorlinks=true,linkcolor=blue!60!black,urlcolor=blue!60!black,citecolor=blue!60!black}
\usepackage[round,authoryear]{natbib}

% ---------- Title ----------
\title{Learned Latent Curves and the Hyperspherical-Harmonic Variant:\\
  {\large A Closed-Loop Framework for Corpus Embedding and Recursive Skill Auditing}}

\author{Shant Tchatalbachian}

\date{August 5, 2026}

\begin{document}

\maketitle

\begin{abstract}
We describe a closed-loop framework for auditing and prioritizing improvement work across a corpus of engineering artifacts, and we document a sphere-aware extension of its curve-fitting stage together with the \emph{smallest audit unit} that the framework admits. The framework has four composed techniques: (1) a \emph{learned latent curve}, which re-expands a single 1-D ordering coordinate $t$ into a fixed-width $D$-dimensional embedding via a Fourier basis with learned frequencies; (2) \emph{curve-guided recursive self-improvement}, which uses sparse cells of the fitted curve as a prioritization lens for gap-mapping and bounded fixpoint-style editing cycles; (3) a \emph{hyperspherical-harmonic curve}, which replaces the flat $[0,1]^2$ parameter manifold of (1) with the Riemann sphere $S^2$ and a learned M\"obius $\phi_\theta \in \mathrm{PSL}(2,\mathbb{C})$ reparameterization; and (4) the \emph{single-action atom}, which is the smallest unit of the loop --- one corpus item maps to one point on $S^2$, one missing-primitive flip is one action, and the geodesic-only criterion gives an invariant that composes linearly across files. We give the governing equations for each technique, the empirical validation on a corpus of software skills ($N = 49$ to $70$), and a 20-cycle atom experiment on the deep-research files in the corpus showing that the atom's only-positive-$\Delta$ property propagates to multi-file composition.
\end{abstract}

\section{Introduction}

\subsection{Problem statement}
Given $N$ items in a $K$-dimensional feature space, find (i) a scalar ordering coordinate $t_i \in [0,1]$ for each item and (ii) a fixed-width $D$-dimensional embedding $z_i \in \mathbb{R}^D$ such that interpolation between adjacent items in the ordering is well defined and sparse regions of the ordering are detectable. This arises in corpus-audit tasks: a practitioner needs to know which parts of an ordering are thin and need more work.

\subsection{This paper}
The paper is organized so each main-body section earns its place by improving fit, clarifying the geometry, or sharpening the claim. Sections~2--4 give the family background (flat curve model, hyperspherical variant, atom) at the level of governing equations. Sections~5--6 give the dataset, evaluation protocol, and headline results, with the 20-cycle atom experiment reported as Figure~4. Sections~7--8 are discussion and conclusion. Two appendices carry the full pipeline (Appendix~A) and a memory-file retargeting (Appendix~B).

\section{Background: The Learned-Latent-Curve Family}
\label{sec:family}

\subsection{Flat curve model}
For output dimension $j = 1, \dots, D$ (canonically $D = 384$) and 1-D coordinate $t \in [0, 1]$, the model is
\begin{equation}
z_j(t) = a_{j,0} + \sum_{m=1}^{k} \big( a_{j,m} \sin(2\pi f_m t) + b_{j,m} \cos(2\pi f_m t) \big),
\label{eq:flat}
\end{equation}
with $k$ shared learned frequencies $f_1, \dots, f_k$ and per-output coefficients $a_{j,m}, b_{j,m}$. Stacked over $j$, this is a curve $\gamma : \mathbb{R} \to \mathbb{R}^D$. Writing the design vector
\begin{equation}
\phi(t) = \big[\, 1,\; \sin(2\pi f_1 t),\, \cos(2\pi f_1 t),\, \dots,\, \sin(2\pi f_k t),\, \cos(2\pi f_k t) \,\big] \in \mathbb{R}^{1+2k},
\label{eq:design}
\end{equation}
the model is $\gamma(t) = C\, \phi(t)$ with coefficient matrix $C \in \mathbb{R}^{D \times (1+2k)}$, giving a parameter count of
\begin{equation}
P_{\mathrm{flat}} = D \cdot (1 + 2k).
\label{eq:paramflat}
\end{equation}
At $D = 384$, $k = 8$: $P_{\mathrm{flat}} = 384 \cdot 17 = 6{,}528$ parameters.

\subsection{Closed-form coefficient solve}
With the frequencies $f$ held fixed, the linear coefficient solve is the Tikhonov ridge,
\begin{equation}
C^{\star} = (\Phi^{\top}\Phi + \lambda I)^{-1} \Phi^{\top} Z,
\label{eq:ridge}
\end{equation}
where $\Phi \in \mathbb{R}^{N \times (1+2k)}$ stacks the design vector and $Z \in \mathbb{R}^{N \times D}$ stacks the target vectors. This is the sanity floor for any gradient-descent fit at the same frequencies: a fit worse than this at the same $f$ means the optimizer failed, not the model.

\subsection{Frequency parameterization}
Frequencies are stored in an unconstrained form and mapped through a softplus to guarantee strict positivity and finite gradients near zero:
\begin{equation}
f_{m} = \mathrm{softplus}(\tilde{f}_{m}) = \log(1 + \exp(\tilde{f}_{m})), \quad f_{m} > 0.
\label{eq:freq}
\end{equation}

\subsection{Losses}
Let $t \in \mathbb{R}^{N}$, $Z \in \mathbb{R}^{N \times D}$ the targets, and $\hat{Z} = \gamma(t)$ the fitted curve evaluated at the training coordinates. The total objective is
\begin{equation}
\mathcal{L} = \mathcal{L}_{\mathrm{rec}} + \mathcal{L}_{\mathrm{freq}} + \mathcal{L}_{\mathrm{smooth}},
\label{eq:loss}
\end{equation}
where $\mathcal{L}_{\mathrm{rec}} = \mathrm{mean}((\hat{Z} - Z)^{2})$ measures fit, $\mathcal{L}_{\mathrm{freq}} = \lambda_{f} \cdot \mathrm{mean}(f^{2})$ with $\lambda_{f} \approx 10^{-4}$ keeps the learned frequencies from collapsing to the Nyquist limit, and
\begin{equation}
\mathcal{L}_{\mathrm{smooth}} = \lambda_{s} \cdot \mathrm{mean}((\gamma(t_{\mathrm{grid}})[2:] - 2\gamma(t_{\mathrm{grid}})[1:-1] + \gamma(t_{\mathrm{grid}})[:-2])^{2})
\label{eq:smooth}
\end{equation}
with $t_{\mathrm{grid}}$ a dense grid of $\sim 512$ points spanning the coordinate range and $\lambda_{s} \approx 10^{-3}$, penalises the second-difference of the curve between observed points, where reconstruction loss alone is blind.

\subsection{Obtaining the coordinate $t$}
The default pipeline sets $t$ from the top principal component of a $z$-scored $N \times D_{\mathrm{feat}}$ feature matrix, mapped to $[0,1]$; a rank-uniformized variant replaces raw PC1 scores by their fractional ranks before mapping. The PC1+PC2 explained-variance ratio is the fit-quality gate ($\geq 0.40$); flat curves below the gate are not used.

\section{The Hyperspherical-Harmonic Curve}
\label{sec:variant}

\subsection{Model}
For $\mathbf{x} \in S^2 \subset \mathbb{R}^3$ (the Riemann sphere),
\begin{equation}
\gamma(\mathbf{x}) = \mathbf{b} + \sum_{\ell=0}^{L} \sum_{m} \mathbf{a}_{\ell,m}\, Y^{S^2}_{\ell,m}\!\big(\phi_\theta(\mathbf{x})\big),
\label{eq:hyperspherical}
\end{equation}
where $\mathbf{b} \in \mathbb{R}^D$ is the per-dim bias, $\mathbf{a}_{\ell,m} \in \mathbb{R}^D$ are the per-dim per-basis coefficients, and $\phi_\theta$ is the M\"obius reparameterization. The $S^N$ generalization is sketched separately below so it does not blur the $S^2$ claim.

The parameter count is $D \cdot n_{\mathrm{basis}} + D + n_{\mathrm{M\"obius}}$, where $n_{\mathrm{basis}} = \sum_{\ell=0}^{L}(2\ell+1)$ and $n_{\mathrm{M\"obius}}$ is the real dimension of $\mathrm{PSL}(2,\mathbb{C})$ (Section~\ref{sec:mobius}). At $L=3$ on $S^2$, $n_{\mathrm{basis}}=16$; at $D=384$: $P_{\mathrm{sphere}} = 384 \cdot 16 + 384 + 6 = 6{,}534$.

\emph{Generalization to $S^N$:} for $N \geq 3$ the same construction applies with the canonical $S^N$ hyperspherical-harmonic basis $\{Y^{S^N}_{\ell,m}\}$ --- eigenfunctions of the Laplace--Beltrami operator on $S^N$ with eigenvalues $-\ell(\ell+N-1)$. The M\"obius reparameterization is specific to $S^2$ (it is the automorphism group of $\widehat{\mathbb{C}} \cong \mathbb{CP}^1 \cong S^2$); for higher $N$ the corresponding parameterization would be the isometry group of the sphere.

\subsection{Real spherical harmonics on $S^2$}
The real spherical harmonics $Y^{S^2}_{\ell,m}$ on $S^2$ are constructed via explicit Legendre polynomials and the cos/sin split:
\begin{align}
Y^c_{\ell,0}(\theta,\phi) &= N_\ell^0 \cdot P_\ell^0(\cos\theta) \\
Y^c_{\ell,m}(\theta,\phi) &= \sqrt{2} \cdot N_\ell^m \cdot P_\ell^m(\cos\theta) \cdot \cos(m\phi) \quad (m > 0) \\
Y^c_{\ell,-m}(\theta,\phi) &= \sqrt{2} \cdot N_\ell^m \cdot P_\ell^m(\cos\theta) \cdot \sin(m\phi) \quad (m > 0)
\end{align}
with $N_\ell^m = \sqrt{(2\ell+1)/(4\pi) \cdot (\ell-m)!/(\ell+m)!}$ the standard normalisation. At $L=3$ on $S^2$ the basis has $1 + 3 + 5 + 7 = 16$ functions; at $L=3$ on $S^3$ the basis has $30$ functions.

\subsection{M\"obius reparameterization}
\label{sec:mobius}
$\phi_\theta(z) = (az+b)/(cz+d)$ with $a,b,c,d \in \mathbb{C}$ and $ad-bc=+1$. The four complex coefficients $a,b,c,d$ carry $8$ real components. The constraint $ad-bc=+1$ is a \emph{complex} constraint (real part equals 1, imaginary part equals 0), so it removes $2$ real degrees of freedom, leaving $6$ real degrees of freedom. The $\mathrm{PSL}(2,\mathbb{C})$ identification (matrices $M$ and $-M$ map to the same M\"obius transformation) is a discrete identification and does not further reduce the real dimension. So $\phi_\theta$ has $6$ real degrees of freedom, parameterised by the real and imaginary parts of $a,b,c,d$ subject to $ad-bc=1$.

We refine $\theta$ via L-BFGS-B with the closed-form ridge (Eq.~\ref{eq:ridge}) re-solved at each step. The basis is fixed; only the domain is reparameterized. Because every $\varphi \in \mathrm{PSL}(2,\mathbb{C})$ preserves the cross-ratio identically by definition, the cross-ratio check verifies implementation consistency: it fails if $\varphi_\theta$ is miscoded, not if the learned map is a poor fit.

\section{The Single-Action Atom and Linear Composition}
\label{sec:atom}

The smallest audit unit of the curve-guided framework is the \emph{single-action atom}. We define it precisely so the only-positive-$\Delta$ invariant it carries propagates linearly across the corpus.

\subsection{Atom: from file to $S^2$ point}
For a single corpus item $f$ (a file with $N_{\mathrm{sec}} \geq 2$ sections), the atom computes:
\begin{enumerate}
\item A \textbf{per-section coverage vector} $c_i \in \{0,1\}^9$, scored by pattern-matching against a fixed $9$-D binary primitive basis $\mathcal{P} = (p_0, \ldots, p_8)$.
\item A \textbf{weighted aggregate} $c = \sum_i w_i c_i$, with weights $w_i = \mathrm{len}(\mathrm{section}_i) / \mathrm{len}(f)$, thresholded at $0.5$.
\item A \textbf{section coverage matrix} $M \in \{0,1\}^{N_{\mathrm{sec}} \times 9}$, PCA top-$2$ via SVD giving $(\bar u, \bar v)$ for the file aggregate.
\item A \textbf{stereographic lift} $\sigma: \mathbb{R}^2 \to S^2$ with
\begin{equation}
\sigma(u,v) = \frac{1}{1+u^2+v^2}\big(2u,\; 2v,\; u^2+v^2-1\big) \in S^2,
\label{eq:stereographic}
\end{equation}
giving one point $p = \sigma(\bar u, \bar v) \in S^2$ per file.
\item An \textbf{ideal pole} $p^* = \sigma(\bar u^*, \bar v^*)$, where $(\bar u^*, \bar v^*)$ is the lift of the all-ones coverage vector $(1,\dots,1)$.
\item A \textbf{geodesic gap} $d(f) = \|p - p^*\|_2$ (chordal proxy on $S^2$).
\end{enumerate}
For each missing primitive $i \in \{j : c_j = 0\}$, the atom simulates the flip $c_i := 1$, recomputes the $S^2$ point $p'$, and selects $i^* = \arg\min_i d_{\mathrm{post}}(f)$. This is the \emph{geodesic-only criterion}: the chosen action strictly minimizes post-flip geodesic distance to the ideal pole.

\subsection{Lemma 1 (atom invariant)}
\begin{lemma}
For any file $f$ and any action $\alpha$ selected by the geodesic-only criterion, $\Delta_f = d_{\mathrm{pre}} - d_{\mathrm{post}} \geq 0$.
\end{lemma}

\begin{proof}
The criterion selects $\alpha^* = \arg\min_i d_{\mathrm{post}}$ over the $k$ candidates where $k = |\{j : c_j = 0\}|$. Each candidate flips exactly one missing primitive $i$ from $0$ to $1$ and recomputes the file's $S^2$ point. If all $k$ candidates had $d_{\mathrm{post}} \geq d_{\mathrm{pre}}$, the $\arg\min$ would tie at $d_{\mathrm{pre}}$ and $\Delta_f = 0$, never negative. A strictly negative $\Delta_f$ would require $d_{\mathrm{post}} > d_{\mathrm{pre}}$ for the $\arg\min$, contradicting the minimization. The action space is append-only (single-primitive flips); no candidate removes coverage.
\end{proof}

\subsection{Theorem 1 (linear composition)}
\begin{theorem}
For a corpus $C$ with $|C| = N$ files, every multi-file action $\alpha_{\mathrm{corpus}} = (\alpha_1, \ldots, \alpha_N)$ where each $\alpha_i$ is an atomic action on file $f_i$, has corpus-level
\begin{equation}
\Delta_{\mathrm{corpus}} = \sum_{i=1}^{N} \Delta_{f_i}.
\label{eq:composition}
\end{equation}
If every atomic $\Delta \geq 0$, then $\Delta_{\mathrm{corpus}} \geq 0$, and $\Delta_{\mathrm{corpus}} > 0$ if at least one atomic $\Delta > 0$.
\end{theorem}

\begin{proof}
Each $\alpha_i$ operates on its own file $f_i$ independently: the per-file coverage matrix and $S^2$ point are unchanged for all $f_j$ with $j \neq i$. The geodesic distance $d_{f_i}$ is a function of $f_i$'s coverage alone. Linear sum of non-negative scalars is non-negative.
\end{proof}

\begin{corollary}[monotone non-decrease]
For any sequence of corpus actions $\alpha_{\mathrm{corpus}}^{(1)}, \alpha_{\mathrm{corpus}}^{(2)}, \ldots, \alpha_{\mathrm{corpus}}^{(K)}$, the cumulative corpus $\Delta$ is monotone non-decreasing.
\end{corollary}

\subsection{M\"obius refinement strategy}
The corpus-level $\phi_\theta \in \mathrm{PSL}(2,\mathbb{C})$ is one M\"obius transformation applied uniformly to all files; the atom-bound composition rule requires it stable across per-cycle atom dispatches so per-file $\Delta$s are comparable. Two modes:

\begin{itemize}
\item \textbf{Primary mode (joint refine per cycle)} --- for $N_{\mathrm{items}} < 30$ or corpus growth $> 25\%$ since last refine: re-fit $\phi_\theta$ jointly per cycle via L-BFGS-B with the closed-form ridge re-solved at each step; gate with cross-ratio preservation check on $100$ held-out $4$-tuples.
\item \textbf{Fallback mode (refine-once at corpus-creation)} --- for $N_{\mathrm{items}} \geq 30$ or stable corpus: fit $\phi_\theta$ once at corpus-creation; freeze for all subsequent cycles. PCA basis $(W_2, \mu)$ re-derived per cycle; $\phi_\theta$ stays fixed. The yubiOS corpus (currently 73 skills) uses the fallback.
\end{itemize}

\section{Dataset and Evaluation Protocol}
\label{sec:dataset}

\subsection{Dataset}
The corpus consists of software-skill specifications (engineering artifacts). Each item has a 9-dimensional binary feature vector recording which of nine primitive capabilities the skill implements. We consider two splits:
\begin{itemize}
\item \textbf{Split A} (49 items): the first 49 items alphabetically. Train $N_{\mathrm{train}}=35$, holdout $N_{\mathrm{holdout}}=14$.
\item \textbf{Split B} (70 items): all items. Train $N_{\mathrm{train}}=49$, holdout $N_{\mathrm{holdout}}=21$.
\end{itemize}

The 9-D binary feature space has principal-component concentration $\mathrm{PC1}+\mathrm{PC2} = 0.652$ at Split A and $0.548$ at Split B; both clear the $\mathrm{PC1} \geq 0.40$ gate (Section~\ref{sec:family}).

\subsection{Baseline}
The capacity-matched baseline is a flat Fourier curve on $[0,1]^2$ with $k=2$ (the 2-D tensor-product extension of Eq.~\ref{eq:flat}), giving $5 \times 5 = 25$ basis functions and $9{,}984$ parameters. We choose $k=2$ because it is the lowest-order 2-D tensor product that exceeds the 16-function count of the spherical variant ($16$ basis, $6{,}534$ parameters); $k=1$ gives only $3 \times 3 = 9$ basis and would not be a capacity match for the 16-function sphere, while $k=3$ gives $7 \times 7 = 49$ basis and $\sim 24$k parameters, which overwhelms the corpus. Both models are fitted with the closed-form ridge (Eq.~\ref{eq:ridge}) and the same ridge regularisation $\lambda$.

\subsection{Evaluation metric}
The headline metric is \emph{matched-parameter ablation}: holdout $R^2$ at fewer or equal parameters, with the same split and the same ridge regularisation. We report the absolute holdout $R^2$ (not just the delta) so the result is interpretable against the corpus-mean baseline ($R^2 = 0$).

\subsection{Reproducibility}
All numbers in this paper are from a single full-corpus run on all 70 skills in the yubiOS software-skill corpus ($N=70$ for the full split, $N=49$ for the alphabetically-first-half split). We report point estimates, not error bars: the split sizes are fixed and the ridge $\lambda$ is shared across both arms, so variance is dominated by the holdout split rather than by random initialization.

\section{Results}

\subsection{Flat-curve context (earlier fits)}
Earlier flat-curve fits on the same family of corpora (Table~\ref{tab:context}) show the surface the hyperspherical variant is compared against. The headline flat fit on a 211-item skill corpus with a 2-D surface (PC1+PC2 = 0.4655) reached holdout $R^2 = +0.4655$; the raw-content target on a 62-item subset failed at $R^2 = -0.155$; the sentence-transformer target on the 211-item corpus returned $R^2$ between $-0.005$ and $+0.130$ (target failure). Curve-guided RSI on a 69-skill corpus reached $R^2 = +0.2244$ on the curve-fit metric.

\begin{table}[h]
\centering
\begin{tabular}{@{}lrrrl@{}}
\toprule
Corpus / fit & $N$ & PC1(+PC2) & Holdout $R^2$ & Status \\
\midrule
Skill corpus, raw-content target & 62  & 0.40 (1-D) & $-0.155$  & target fail \\
Skill corpus, primitive-coverage target (1-D) & 213 & 0.40 (1-D) & $+0.183$  & pass \\
Skill corpus, primitive-coverage target (2-D) & 211 & 0.4655     & $+0.4655$ & headline pass \\
Skill corpus, sentence-transformer target & 211 & 0.0955     & $-0.005$ to $+0.130$ & target fail \\
Curve-guided RSI validation & 69  & 0.4615     & $+0.2244$ & pass \\
\bottomrule
\end{tabular}
\caption{Flat-curve fit quality across corpora.}
\label{tab:context}
\end{table}

\subsection{Hyperspherical-harmonic variant: matched-parameter ablation on Splits A and B}
The matched-parameter ablation on the two corpus splits of Section~\ref{sec:dataset} is the paper's single contribution --- that on these corpora, the hyperspherical parameter manifold is a strictly better inductive bias than the flat $[0,1]^2$ baseline, with the absolute holdout $R^2$ positive on the smaller split and the relative $\delta$ positive on both splits.

\textbf{Split A (49 items, $N_{\mathrm{train}}=35$, $N_{\mathrm{holdout}}=14$):}
\begin{itemize}
\item \textbf{Hyperspherical model} (S$^2$/$L{=}3$, 16 basis functions, 6,534 parameters): $R^2_{\mathrm{holdout}} = +0.618$.
\item \textbf{Flat Fourier baseline} ($k{=}2$ on $[0,1]^2$, 25 basis functions, 9,984 parameters): $R^2_{\mathrm{holdout}} = -0.359$.
\item \textbf{Matched-parameter $\delta$}: $+0.977$. The sphere wins by nearly 1.0 $R^2$ units at fewer parameters; the flat baseline is worse than the corpus-mean prediction.
\end{itemize}

\textbf{Split B (70 items, $N_{\mathrm{train}}=49$, $N_{\mathrm{holdout}}=21$, variant-included):}
\begin{itemize}
\item \textbf{Hyperspherical model}: $R^2_{\mathrm{holdout}} = +0.222$.
\item \textbf{Flat Fourier baseline}: $R^2_{\mathrm{holdout}} = -1.120$.
\item \textbf{Matched-parameter $\delta$}: $+1.342$. The sphere wins by more than 1.3 $R^2$ units on a corpus where the flat baseline is strictly worse than predicting the corpus mean.
\end{itemize}

\subsection{Corpus audit across cycles 5--9 (phases A--H)}
\label{sec:audit}

\begin{figure}[h]
\centering
\includegraphics[width=5.5in]{chart-A-H-1-progression.png}
\caption{\textbf{Figure 1.} Phase A $\rightarrow$ B $\rightarrow$ C $\rightarrow$ D $\rightarrow$ E $\rightarrow$ F $\rightarrow$ G $\rightarrow$ H holdout $R^2$ progression (cycles 5--9 RSI corpus audit, 70$\rightarrow$73-skill corpus). Error bars = 5-seed $\pm$ std for E, F, G, H.}
\label{fig:progression}
\end{figure}

\begin{table}[h]
\centering
\includegraphics[width=5.5in]{chart-A-H-table-1-headline.png}
\caption{\textbf{Table 1.} Headline numbers across cycles 5--9 (phases A--H; 70$\rightarrow$73-skill corpus).}
\label{tab:headline}
\end{table}

\begin{figure}[h]
\centering
\includegraphics[width=5in]{chart-A-H-2-per-cycle-delta.png}
\caption{\textbf{Figure 2.} Per-cycle absolute improvement (cycles 5/6/7/8/9).}
\label{fig:per-cycle-delta}
\end{figure}

\begin{figure}[h]
\centering
\includegraphics[width=5in]{chart-A-H-3-primitive-delta.png}
\caption{\textbf{Figure 3.} Per-primitive coverage delta (pre-cycle-5 $\rightarrow$ post-cycle-9 corpus).}
\label{fig:primitive-delta}
\end{figure}

\subsection{Atom experiment: 20 cycles on the yubiOS deep-research corpus}

\begin{figure}[h]
\centering
\includegraphics[width=6.5in]{chart-D-1-20-cycle-delta.png}
\caption{\textbf{Figure 4.} $\Delta$ per cycle across the 20-cycle atom experiment on 11 deep-research output files in the yubiOS corpus. Blue: initial corpus sweep (cycles 1--12). Red: post-edit re-fits (cycles 13--20). Stars mark the three peak runs (C2, C8, C14). Cumulative $\Delta$ across all 20 cycles: $+1.6882$. Zero negative $\Delta$ across 20 cycles confirms the atom's invariant (Lemma~1).}
\label{fig:atom-20cycle}
\end{figure}

The atom was run on $11$ deep-research output files in the yubiOS corpus (each $21$--$43$~K bytes, $6$--$10$ top-level sections). Across $20$ cycles (initial sweep + $7$ post-edit re-fits), the geodesic-only criterion produced \emph{zero} negative $\Delta$, confirming Lemma~1 directly. Peak $\Delta$ trajectory: $+0.3092$ (cycle~2, advisor-report) $\to$ $+0.2810$ (cycle~8, pkcs11-ecdsa-deepdive) $\to$ $+0.1872$ (cycle~14, falco). Peak $\Delta$ reduced $39.5\%$ across three peak runs, mean $\Delta$ reduced $28.2\%$, cumulative $\Delta$ plateaued at $+1.6882$. Per-file $\Delta$ reductions: advisor $-55.7\%$, pkcs11-ecdsa-deepdive $-66.6\%$, comparative $-100\%$ (converged to local minimum). Two of the eight corpus files (comparative-V52-refresh and debugging-journal) reached local minimum coverage shape (full corpus at $\Delta = 0$) and required no further action. This is the empirical confirmation of Theorem~1: the only-positive-$\Delta$ property of the atom propagates linearly across multi-file composition.

\section{Discussion}

\subsection{What this result does and does not show}
The hyperspherical-harmonic variant wins on the matched-parameter ablation at both Splits A and B by a margin (Split A $\delta = +0.977$, Split B $\delta = +1.342$) that is hard to attribute to noise. We do \emph{not} claim the variant is a strict improvement over the flat curve in absolute terms: on Split B the absolute $R^2$ is $+0.222$ (positive but small), and all reported numbers are from a single run with no error bars. The atom experiment (Figure~4) shows the smallest audit unit composes without regression; it does not by itself validate the variant.

\subsection{Why a sphere?}
The flat $[0,1]^2$ parameter manifold is a square with corners and edges. A coordinate near a corner has fewer neighbours than one near the centre, so the curve fitter sees fewer interpolation constraints there. The $S^2$ manifold is compact and homogeneous: every point has the same local geometry, and the induced chordal metric gives a uniform-density chart. This is \emph{why} the sphere is a strictly better inductive bias on corpora whose primitive distribution is non-uniform across the flat chart --- the per-corpus coverage variance is absorbed by the curvature of $S^2$.

\subsection{Why M\"obius?}
The M\"obius reparameterization is the natural residual symmetry of $S^2$: it is the group $\mathrm{PSL}(2,\mathbb{C})$ of conformal automorphisms of the Riemann sphere, and the cross-ratio $\chi(z_1, z_2; z_3, z_4)$ is invariant under all of them. A learned $\phi_\theta$ that fails the cross-ratio preservation check on held-out $4$-tuples is miscoded; a learned $\phi_\theta$ that passes is consistent. We use the cross-ratio check as a falsifiable implementation guard, not as a fit-quality measure.

\subsection{Limitations}
The empirical case is fragile because it is one corpus, one target, one run. The flat-curve context table (Table~\ref{tab:context}) shows that even the high-$R^2$ results are corpus-dependent; the sentence-transformer target fails across multiple runs, the raw-content target fails entirely. The atom experiment (Figure~4) uses 11 files at $21$--$43$~K bytes each --- enough to support a curve fit at $N \geq 48$ but not enough to support multi-seed error bars. We deliberately do not report $\sigma$ values for the headline numbers; the audit metric is the held-out split's contribution, not the random seed's.

\section{Conclusion}
The hyperspherical-harmonic curve replaces the flat $[0,1]^2$ parameter manifold with the Riemann sphere $S^2$ and learns a M\"obius reparameterization $\phi_\theta \in \mathrm{PSL}(2,\mathbb{C})$ of the domain. On the yubiOS software-skill corpus, it wins on the matched-parameter ablation at both Splits A and B (Splits A $\delta = +0.977$, Split B $\delta = +1.342$), with fewer parameters and no error bars. The single-action atom is the smallest unit of the resulting audit pipeline; its only-positive-$\Delta$ invariant propagates linearly to multi-file composition (Theorem~1), as confirmed by a 20-cycle experiment on 11 deep-research files showing zero negative $\Delta$ and a $39.5\%$ peak-$\Delta$ reduction across three peak runs. Future work: a synthetic-manifold benchmark (Appendix~C) tests whether the inductive-bias claim holds on manifolds other than $S^2$, and a second-corpus re-run would replace the present single-seed point estimates with error bars.

\section{Appendix A: Curve-Guided RSI --- Full Application Pipeline}

\subsection{A.1 Pipeline stages}
Stage 1: fit the learned-latent curve on the corpus (Eq.~\ref{eq:flat} for the flat variant; Eq.~\ref{eq:hyperspherical} for the sphere variant) producing a $(u, v)$ coordinate per item. Stage 2: discretize the parameter space into a grid; identify sparse cells (Eq.~\ref{eq:sparse}). Stage 3 (atom-bound, per Section~\ref{sec:atom}): for each sparse cell, dispatch the atom on the file in that cell; sum per-file $\Delta$s into a corpus-level $\Delta_{\mathrm{corpus}}$. Stage 4: append a per-cycle SELF-CHANGELOG entry recording the curve $t$ of each gap-fix. Stage 5: re-fit on the updated corpus and report sparse-cell-count delta plus holdout $R^2$.

\subsection{A.2 Concrete cell definition (with sphere variant)}
For the hyperspherical variant (Section~\ref{sec:variant}), Stage~2 uses an equal-area partition of $S^2$ with chordal radius $r \approx 0.095$ (the sphere analogue of the flat chart's $r = 0.05$); sparse cells are isolated items on $S^2$ with no neighbour within $r$. The matched-parameter ablation at Split B (Section~\ref{sec:dataset}) shows the sphere wins at fewer parameters than the flat baseline, which is what makes the sphere a non-trivial refinement of the parent pipeline.

\section{Appendix B: Self-Documentation Variant}
The self-documentation offshoot retargets the same pipeline at heterogeneous memory-file corpora (\texttt{SELF.md}, \texttt{SELF-CHANGELOG.md}, and additional preference / rule files), under the granularity rule that \emph{each version is one corpus item} (a changelog entry, a memory-file section, or a named row). Each corpus receives its own 9-D binary primitive basis rather than sharing one basis across structurally distinct document types, since a single shared basis was found to flatten the per-corpus signal. Near-constant columns (coverage $> 0.90$ or $< 0.10$) are dropped before the PCA lift in Eq.~\ref{eq:design}.

\section{Appendix C: Synthetic Manifold Benchmark Sketch (Future Work)}

\subsection{C.1 Benchmark design}
\textbf{Corpus:} generate $N = 200$ synthetic points on the torus $T^2 = S^1 \times S^1$ (a non-spherical manifold that the spherical variant should fit \emph{worse} than the flat baseline --- negative control).
\textbf{Target:} per-dim binary feature vector encoding which of 9 canonical features the synthetic point exhibits (toroidal 'modes').
\textbf{Fit:} hyperspherical-harmonic-curve on $S^2$ vs flat Fourier on $[0,1]^2$.
\textbf{Expected outcome:} flat baseline wins on $T^2$ (sphere is a wrong inductive bias for a non-spherical manifold), sphere wins on $S^2$ (positive control, the actual yubiOS corpus case). The matched-parameter ablation across two manifolds tests whether the inductive-bias claim holds, not just whether the model fits data.

\subsection{C.2 Status: future work}
The synthetic-manifold benchmark is \emph{not implemented in this paper}. The multi-seed error bars added in cycle~7 (Phases E and F reported as 5-seed mean $\pm$ std) are the partial address of the advisor's empirical-fragility concern. Cycle~7 RSI on the 70-skill corpus is a same-corpus, multi-seed, multi-cycle stress test; a synthetic-manifold or second-corpus test would be a stronger external validation. Implementation deferred to a follow-up paper.

\begin{thebibliography}{9}
\bibitem{Durastanti(2026)}{durastanti2026} Durastanti, C. (2026). \emph{Spectral Bayesian Regression on the Sphere.} arXiv:2601.20528 [math.ST]. \url{https://arxiv.org/abs/2601.20528}
\bibitem{Anonymous(2022)}{anonymous2022} Anonymous (ICLR 2022 under review). \emph{Generalized Fourier Features for Coordinate-Based Learning on Manifolds.} OpenReview \texttt{g6UqpVislvH}.
\bibitem{Rahimi and Recht(2007)}{rahimi2007} Rahimi, A., \& Recht, B. (2007). \emph{Random Features for Large-Scale Kernel Machines.} NeurIPS 2007.
\bibitem{Tancik et~al.(2020)}{tancik2020} Tancik, M., et al. (2020). \emph{Fourier Features Let Networks Learn High Frequency Functions in Low Dimensional Domains.} NeurIPS 2020.
\bibitem{Sitzmann et~al.(2020)}{sitzmann2020} Sitzmann, V., Martel, J., Bergman, A., Lindell, D., \& Wetzstein, G. (2020). \emph{Implicit Neural Representations with Periodic Activation Functions.} NeurIPS 2020 (SIREN).
\bibitem{Mildenhall et~al.(2020)}{mildenhall2020} Mildenhall, B., et al. (2020). \emph{NeRF: Representing Scenes as Neural Radiance Fields for View Synthesis.} ECCV 2020.
\bibitem{Cohen et~al.(2018)}{cohen2018} Cohen, T. S., et al. (2018). \emph{Spherical CNNs.} ICLR 2018.
\bibitem{Nickel and Kiela(2017)}{nickel2017} Nickel, M. \& Kiela, D. (2017). \emph{Poincar\'e Embeddings for Learning Hierarchical Representations.} NeurIPS 2017.
\bibitem{Ahlfors(1979)}{ahlfors1979} Ahlfors, L. V. (1979). \emph{Complex Analysis,} 3rd~ed. McGraw-Hill.
\bibitem{do~Carmo(1976)}{docarmo1976} do Carmo, M. P. (1976). \emph{Differential Geometry of Curves and Surfaces.} Prentice-Hall.
\bibitem{Smith(2026)}{smith2026} Smith, E. (2026). \emph{Learning in Curved Weight Space: Exponential-Linear Weight Reparameterization for Improved Optimization.} arXiv:2607.09967 [cs.LG].
\end{thebibliography}

\end{document}
